\chapter{Backend Implementation}
\label{ch:backend}

This chapter documents each backend service in depth. For every implemented
service it describes the responsibility, technology stack, data model, REST and
WebSocket API, event interactions, and the notable business logic that gives the
service its character. The two unimplemented scaffolds and the placeholder shared
library are described honestly for completeness. Throughout, the cross-cutting
patterns of Chapter~\ref{ch:architecture} --- gateway-injected identity,
ticket-based sockets, the \code{\{action, payload\}} envelope, and idempotent
event handling --- appear in their concrete form.

\section{Gateway API}
\label{sec:be-gateway}

The gateway (\code{gateway-api}, port 8080) is a reactive Spring Cloud Gateway
application with no datastore. Its routing table maps public path prefixes to
downstream services (Table~\ref{tab:gw-routes}), each target overridable by an
environment variable for containerised deployment.

\begin{table}[htbp]
\centering
\caption{Gateway routing rules.}
\label{tab:gw-routes}
\small
\begin{tabularx}{\textwidth}{L{3.4cm} X L{3.2cm}}
\toprule
\textbf{Route} & \textbf{Path predicate} & \textbf{Target (default)} \\
\midrule
user-service & \route{/api/auth/**}, \route{/api/otp/**}, \route{/api/users/**} & \code{localhost:8091} \\
workspace-service & \route{/api/workspace/**} & \code{localhost:8087} \\
chat-service & \route{/api/chat/**} & \code{localhost:8084} \\
room-service & \route{/api/rooms/**} & \code{localhost:8086} \\
\bottomrule
\end{tabularx}
\end{table}

The security core is a single global filter, \code{AuthGatewayFilter}, ordered
ahead of routing. Its logic captures the trust model of the whole platform:

\begin{lstlisting}[language=Java,caption={Representative gateway security filter (the platform's trust boundary).}]
// order -1: runs before routing
public Mono<Void> filter(ServerWebExchange exchange, GatewayFilterChain chain) {
    ServerHttpRequest req = exchange.getRequest();
    String path = req.getURI().getPath();

    // 1. anti-spoofing: never let a client set identity headers
    ServerHttpRequest cleaned = req.mutate()
        .headers(h -> { h.remove("X-User-Id"); h.remove("X-User-Role"); })
        .build();

    // 2. internal endpoints are unreachable from the edge
    if (path.startsWith("/api/internal/"))
        return forbidden(exchange);            // 403

    // 3. public passthrough
    if (path.endsWith("/health") || isPublicPrefix(path))
        return chain.filter(exchange.mutate().request(cleaned).build());

    // 4. authenticate: require + verify Bearer JWT, inject trusted identity
    return jwtService.verifyAndExtractUserId(bearer(req))
        .map(userId -> cleaned.mutate()
            .header("X-User-Id", userId.toString())
            .header("X-User-Role", "USER")
            .build())
        .map(r -> chain.filter(exchange.mutate().request(r).build()))
        .orElse(unauthorized(exchange));       // 401
}
\end{lstlisting}

The JWT itself is verified with HS256 using a shared secret (\code{jwt.secret})
identical to the one the user service signs with; the filter reads only the
\code{userId} claim. A \code{CorsWebFilter} runs even earlier so that preflight
\code{OPTIONS} requests never reach the auth filter. Because the gateway hard-codes
the injected role to \code{USER}, platform-level RBAC is not yet modelled; all
role decisions happen per-workspace in the workspace service
(Section~\ref{sec:be-workspace}).

\section{User Service}
\label{sec:be-user}

\paragraph{Responsibility.} The user service (port 8091) is the identity
authority: registration, login and JWT issuance, profile and profile-picture
management, password change, and OTP-based verification. It is the origin of
every JWT in the system and the producer of user-lifecycle notification events.

\paragraph{Stack.} Spring Boot with Web, Data JPA, Security, and AMQP; a MySQL
datastore; \code{jjwt} for tokens; and BCrypt for password and OTP hashing. It
has no WebSocket surface and sends no e-mail itself --- that is delegated to the
notification service via events.

\paragraph{Data model.} Two JPA entities (Table~\ref{tab:user-entities}). A
\code{User} owns many \code{VerificationRequest}s; the profile picture is stored
inline as a \code{LONGBLOB} with its MIME type.

\begin{table}[htbp]
\centering
\caption{User service entities.}
\label{tab:user-entities}
\scriptsize
\begin{tabularx}{\textwidth}{L{2.4cm} X}
\toprule
\textbf{Entity (table)} & \textbf{Fields} \\
\midrule
\code{User} (\code{users}) & \code{id}, \code{firstName}, \code{lastName}, \code{email} (unique), \code{phoneNumber}, \code{password} (BCrypt), \code{accountStatus} (\code{ACTIVE}/\code{INACTIVE}/\code{PENDING\_REPORT\_REVIEW}/\code{SUSPENDED}), \code{isEmailVerified}, \code{isPhoneVerified}, \code{isDisabled}, \code{profilePicture} (\code{byte[]}), \code{profilePictureType} \\
\code{Verification\-Request} (\code{verification\_requests}) & \code{id}, \code{otp} (BCrypt hash), \code{status} (\code{PENDING}/\code{APPROVED}), \code{createdAt}, \code{expiresAt} (+10\,min), \code{type} (\code{PASSWORD\_RESET}/\code{EMAIL\_VERIFICATION}/\code{CHANGE\_EMAIL}), \code{user} (\code{@ManyToOne}) \\
\bottomrule
\end{tabularx}
\end{table}

\paragraph{API.} The REST surface (Table~\ref{tab:user-api}) splits into public
authentication endpoints and protected profile endpoints; the latter are guarded
by the service's own \code{JwtAuthFilter} for defence in depth, in addition to
the gateway.

\begin{table}[htbp]
\centering
\caption{User service REST API.}
\label{tab:user-api}
\scriptsize
\begin{tabularx}{\textwidth}{L{1.2cm} L{3.6cm} X}
\toprule
\textbf{Method} & \textbf{Path} & \textbf{Purpose} \\
\midrule
POST & \route{/api/auth/register} & Register, hash password, issue JWT, emit \code{SIGNUP\_SUCCESS}; \code{409} if e-mail exists \\
POST & \route{/api/auth/login} & Authenticate and issue JWT; \code{401} on failure \\
POST & \route{/api/auth/otp/request} & Generate and e-mail an OTP; emit \code{OTP} \\
POST & \route{/api/auth/otp/verify} & Verify an OTP, mark \code{APPROVED} \\
PUT & \route{/api/users/me/password} & Change password (checks old); emit \code{PASSWORD\_RESET\_SUCCESS} \\
GET & \route{/api/users/me} & Current profile \\
POST/GET & \route{/api/users/me/photo} & Upload / return profile picture ($\le$10\,MB) \\
GET & \route{/api/users} & List users (id, name, e-mail) \\
\bottomrule
\end{tabularx}
\end{table}

\paragraph{Authentication internals.} On login, after a BCrypt password check,
\code{JwtUtil} mints an HS256 token whose subject is the e-mail, with a custom
\code{userId} claim and a 24-hour expiry. This is the token the gateway later
verifies. Security is stateless (no server session); a
\code{CustomUserDetailsService} loads users by e-mail and rejects disabled
accounts, mapping a \code{SUSPENDED} status onto a locked account. A
\code{GlobalExceptionHandler} maps authentication failures to \code{401},
integrity violations (duplicate e-mail) to \code{409}, and validation errors to
\code{400}.

\paragraph{Events produced.} On the three lifecycle actions, the service
publishes a \code{NotificationEvent} to the \code{notifications.exchange}:

\begin{lstlisting}[language=json,caption={A \code{NotificationEvent} (OTP request).}]
{ "eventId": "9f1c...-uuid", "type": "OTP", "occurredAt": "2026-07-10T12:00:00Z",
  "payload": { "email": "user@x.com", "firstName": "Sara",
               "otp": "418306", "expiresInMinutes": 10 } }
\end{lstlisting}

\paragraph{Housekeeping.} A scheduled task runs hourly to bulk-delete expired
verification requests, keeping the table small. Configuration points MySQL at
\code{virtual\_office}, uses JPA \code{ddl-auto=update} (a convenience with a
known schema-drift risk in production), and connects to RabbitMQ on the default
port.

\section{Workspace Service}
\label{sec:be-workspace}

\paragraph{Responsibility.} The workspace service (port 8087) is the spatial and
organisational backbone of the platform. It owns \emph{workspaces} (companies),
\emph{desks} (a user's seat, profile, position, and workspace role),
\emph{teams}, \emph{invitations}, and the editable \emph{2D map layout}
(tilesets, layers, zones, spawn points, map objects). It is also the authority on
\emph{workspace roles}, and it exposes an internal API that boots and
synchronises the Colyseus game sessions.

\paragraph{Stack.} The most conventionally ``enterprise'' service: Spring Web,
Data JPA, Validation, AOP, and AMQP over \textbf{PostgreSQL}, with
\textbf{Flyway} owning the schema (\code{ddl-auto=validate}), \textbf{MapStruct}
for DTO mapping, and \textbf{springdoc-openapi} publishing a Swagger UI. Its
build enforces 90\% test coverage via JaCoCo.

\paragraph{Data model.} The richest schema in the system
(Table~\ref{tab:ws-entities}); Figure~\ref{fig:ws-er} sketches the central
relationships. The \code{Workspace} carries an optimistic-locking
\code{layoutVersion}; the \code{Desk} is the join between a user and a workspace
and holds the workspace-scoped \code{role}.

\begin{table}[htbp]
\centering
\caption{Principal workspace-service entities (PostgreSQL).}
\label{tab:ws-entities}
\scriptsize
\begin{tabularx}{\textwidth}{L{2.2cm} X}
\toprule
\textbf{Entity} & \textbf{Key fields} \\
\midrule
\code{Workspace} & \code{id}, \code{name}, \code{slug} (unique), \code{ownerId}, \code{status}, \code{visibility} (\code{INVITE\_ONLY}), \code{inviteToken}, \code{tileSize}, \code{mapWidth/Height}, \code{mapGeometry} (jsonb), \code{layoutVersion} (\code{@Version}) \\
\code{Desk} & \code{userId}, \code{workspaceId}, \code{fullName}, \code{avatarCharacter}, \code{status}, \code{statusEmoji}, \code{positionX/Y}, \code{isOnline}, \code{lastSeenAt}, \code{role} (\code{GUEST}$<$\code{MEMBER}$<$\code{ADMIN}$<$\code{OWNER}), \code{teamId}, \code{inviteStatus}, \code{isActive} --- unique on \code{(userId, workspaceId)} \\
\code{Team} & \code{id}, \code{workspaceId}, \code{name} --- unique on \code{(workspaceId, name)} \\
\code{WorkspaceInvitation} & \code{token} (UUID), \code{email}, \code{status}, 7-day expiry --- one \code{PENDING} per \code{(workspaceId, email)} \\
\code{Zone} & \code{type}, \code{x}, \code{y}, \code{width}, \code{height}, \code{voiceRoomId} (unique), \code{proximityRadius} --- read by the room service \\
\code{MapObject} & \code{type} (\code{COMPUTER}/\code{WHITEBOARD}), \code{positionX/Y}, \code{roomId} (stable UUID), \code{capacity}, \code{isActive} \\
\code{MapLayer}, \code{Tileset}, \code{SpawnPoint} & The Tiled map: gid layer matrices (jsonb), tilesets (unique \code{firstGid}), spawn points \\
\bottomrule
\end{tabularx}
\end{table}

\begin{figure}[htbp]
\centering
\begin{tikzpicture}[font=\scriptsize,
  e/.style={draw=voPrimary,fill=voPrimary!8,rounded corners=2pt,minimum height=0.7cm,
            text width=1.9cm,align=center},
  rel/.style={voGray,-{Stealth[]}}]
\node[e] (u) at (0,2) {User\\(external)};
\node[e] (w) at (3,2.8) {Workspace};
\node[e] (d) at (3,0.9) {Desk};
\node[e] (t) at (6,2.8) {Team};
\node[e] (i) at (6,0.9) {Invitation};
\node[e] (z) at (9,2.8) {Zone};
\node[e] (m) at (9,0.9) {MapObject};
\draw[rel] (d) -- node[above,font=\tiny]{1} (u);
\draw[rel] (d) -- node[right,font=\tiny]{$*$..1} (w);
\draw[rel] (t) -- node[above,font=\tiny]{groups} (d);
\draw[rel] (i) -- node[left,font=\tiny]{to} (w);
\draw[rel] (z) -- (w); \draw[rel] (m) -- (w);
\end{tikzpicture}
\caption{Workspace-service core relationships. A Desk joins a User to a Workspace
and carries the workspace role.}
\label{fig:ws-er}
\end{figure}

\paragraph{API highlights.} Beyond CRUD on workspaces, teams, desks, and map
objects, three areas stand out. \emph{Invitations}: an admin invites by e-mail
(creating a \code{PENDING} invitation and desk and an e-mail event); the invitee
accepts via a tokenised link, which creates or reactivates their desk and fires a
``member added'' event to the chat service. \emph{Layout editing}: replacing the
map is guarded by optimistic concurrency --- the client submits the version it
edited, and a mismatch is rejected --- after which child rows (layers, tilesets,
zones) are validated and replaced atomically. \emph{Internal session API}: under
\route{/api/internal/**} (gateway-blocked, \code{X-Internal-Token}-guarded), the
service serves the Colyseus game server a full \route{/session-config} boot
payload, validates joins, resolves \route{/members/\{userId\}/role} (consumed by
the chat and room services), and ingests avatar \route{/presence} updates.

\paragraph{Events produced.} On workspace creation and membership change, a
\code{WorkspaceChannelEvent} is published --- \emph{after the database
transaction commits} --- to the \code{workspace.exchange}, so the chat service can
maintain a single canonical ``general'' channel per workspace. Publishing
after-commit ensures the chat service never reacts to a change that was later
rolled back, and a broker outage is logged rather than allowed to fail the
user's request.

\section{Chat Service}
\label{sec:be-chat}

\paragraph{Responsibility.} The chat service (port 8084) is the real-time
messaging core: channels (group, direct, and room-bound), messages (with edit,
soft-delete, reply, mention, and idempotent send), threads, and read
receipts/unread counts. It delivers over WebSocket and provisions channels in
response to events from the room and workspace services.

\paragraph{Stack.} Spring Web, WebSocket, Data Mongo, Data Redis, AMQP, and
Validation over \textbf{MongoDB} (documents) and \textbf{Redis} (cursors,
tickets, rate limits). It consumes events but produces none.

\paragraph{Data model.} Three collections (Table~\ref{tab:chat-entities}). A
subtle but important distinction is \code{threadId} (``this message lives inside a
thread'') versus \code{replyToId} (``this message quote-replies another'') --- the
two are independent. Message identity uses MongoDB \code{ObjectId}s, while
channels and threads use integer ids from a sequence-counter collection.

\begin{table}[htbp]
\centering
\caption{Chat service collections.}
\label{tab:chat-entities}
\scriptsize
\begin{tabularx}{\textwidth}{L{1.9cm} X}
\toprule
\textbf{Collection} & \textbf{Fields} \\
\midrule
\code{channels} & \code{id}, \code{name}, \code{type} (\code{GROUP}/\code{DIRECT}/\code{ROOM}), \code{workspaceId}, \code{members} (\code{int[]}), \code{dmKey} (``\code{minId\_maxId}''), \code{canonical}, \code{createdBy}, timestamps \\
\code{messages} & \code{id} (ObjectId), \code{channelId}, \code{senderId}, \code{senderRole} (captured at send), \code{content}, \code{type} (\code{TEXT}/\code{SYSTEM}), \code{threadId}, \code{replyToId}, \code{mentions}, \code{clientMessageId} (idempotency), \code{deleted}, timestamps \\
\code{threads} & \code{id}, \code{channelId}, \code{rootMessageId} (unique), \code{name}, \code{createdBy}, \code{creatorRole}, \code{deleted}, timestamps \\
\bottomrule
\end{tabularx}
\end{table}

Carefully chosen indexes make the store correct and fast: a unique
\code{(senderId, clientMessageId)} index enforces idempotent sends; a unique,
sparse \code{dmKey} index enforces exactly one direct-message channel per user
pair; and a \emph{partial} unique \code{(workspaceId, name)} index that applies
only to \code{GROUP} channels lets room channels reuse names freely.

\paragraph{Authorisation.} Every channel- or thread-scoped operation checks
membership (\code{403} otherwise). Editing is restricted to the author. Deletion
follows a nuanced rule that survives later role changes because the author's role
is \emph{captured on the message}: the owner may always delete; an \code{ADMIN}
may delete others' messages \emph{except} those authored by another \code{ADMIN}.
Before creating a channel or room-bound resource, the service confirms the
caller's workspace role by calling the workspace service's internal role
endpoint, failing closed if it is unreachable.

\paragraph{Read receipts.} Rather than store a read row per message, the service
keeps a single forward-only cursor per user per channel in Redis
(\code{read:\{channelId\}:\{userId\}}), advanced with a Lua script and expiring
after thirty days; the unread count is then simply the number of messages whose
id exceeds the cursor.

\paragraph{Real-time API.} The WebSocket endpoint is \route{/api/chat/connect},
authenticated by the ticket pattern of Section~\ref{sec:realtime-arch}. A
subscription interceptor authorises each \code{SUBSCRIBE} to
\route{/topic/channel/\{id\}} or \route{/topic/thread/\{id\}} against membership,
returning a typed error (\code{NOT\_A\_MEMBER}, \code{CHANNEL\_NOT\_FOUND},
\code{INVALID\_PAYLOAD}) on the user's private \route{/user/queue/errors} stream.
Clients \code{SEND} to \route{/app/chat/send} and \route{/app/chat/typing}; the
server persists and broadcasts, using the uniform envelope with actions
\code{NEW\_MESSAGE}, \code{EDIT\_MESSAGE}, \code{DELETE\_MESSAGE},
\code{THREAD\_DELETED}, and \code{TYPING}. Figure~\ref{fig:chat-send} traces a
send.

\begin{figure}[htbp]
\centering
\begin{tikzpicture}[font=\footnotesize,
  b/.style={draw=voPrimary,fill=voPrimary!8,rounded corners=2pt,minimum height=0.7cm,
            text width=1.9cm,align=center},
  a/.style={-{Stealth[]},thick,voDark},r/.style={-{Stealth[]},densely dashed,voGray},
  note/.style={font=\scriptsize\itshape,voGray}]
\node[b] (s) at (0,0) {Sender};
\node[b] (c) at (3.6,0) {Chat Svc};
\node[b] (m) at (7,0) {MongoDB};
\node[b] (o) at (10.6,0) {Other\\members};
\draw[a] (s) -- node[above,note]{/app/chat/send} (c);
\draw[a] (c) -- node[above,note]{persist} (m);
\node[note] at (3.6,-0.7) {check membership,};
\node[note] at (3.6,-1.0) {idempotent on clientMessageId};
\draw[a] (c) to[out=-30,in=210] node[below,note]{NEW\_MESSAGE @ /topic/channel/id} (o);
\draw[r] (c) to[out=150,in=30] node[above,note]{echo} (s);
\end{tikzpicture}
\caption{Message-send flow in the chat service.}
\label{fig:chat-send}
\end{figure}

\paragraph{Events consumed.} Two listeners keep the channel space in sync with
the rest of the platform. A \code{RoomChannelEvent} (from the room service)
creates, deletes, or edits the membership of a \code{ROOM}-type channel; a
\code{WorkspaceChannelEvent} (from the workspace service) maintains the single
canonical group channel of each workspace. Both are idempotent on \code{eventId}
and retried with backoff, so redelivery is safe and a transient failure
self-heals.

\section{Room Service}
\label{sec:be-room}

\paragraph{Responsibility.} The room service (port 8086) coordinates
voice/video rooms and the proximity-voice layer of the 2D office. It owns room
identity and membership (in MongoDB) and live presence (in Redis), mints Agora
media tokens, pushes presence over WebSocket, computes proximity voice groups
from avatar positions, and publishes room-channel events to the chat service. It
is explicitly a \emph{coordinator}, not a media server --- audio and video never
touch it.

\paragraph{Data model.} A \code{Room} document binds three identities: its own
\code{ObjectId}, the \code{channelId} of the chat channel it is bound to
(generated by the room service and published to chat), and a deterministic
\code{agoraChannelName} (\code{room-<hex>}). Live presence is entirely ephemeral
in Redis: a participants hash per room (each participant's mute, camera, and
screen-share state), a short-lived ``alive'' key per participant refreshed by
heartbeats, and the WebSocket ticket. A capacity limit (default 25) is enforced
\emph{atomically at join} by a Lua script that adds the participant only if the
room is not full.

\begin{lstlisting}[language=Java,caption={Representative atomic capacity check at join (Redis Lua).}]
-- KEYS[1] = room:{id}:participants   ARGV[1] = userId   ARGV[2] = max
if redis.call('HLEN', KEYS[1]) >= tonumber(ARGV[2]) then
    return 0                                   -- ROOM_FULL
end
redis.call('HSET', KEYS[1], ARGV[1], ARGV[3])  -- ARGV[3] = state json
return 1
\end{lstlisting}

\paragraph{API.} Rooms are created, listed, updated, and deleted under
\route{/api/rooms}; members are added and removed; and the central action is
\route{POST /api/rooms/\{id\}/join}, which re-checks workspace membership,
registers presence, \textbf{mints an Agora token}, broadcasts
\code{PARTICIPANT\_JOINED}, and returns the room, the Agora channel name, the
token, and the current participant list. Over WebSocket (\route{/ws/rooms}),
clients send \route{/app/room/state} (mute/camera/screen changes, rebroadcast as
\code{STATE\_CHANGED}) and \route{/app/room/heartbeat} (to refresh their alive
key), and subscribe to \route{/topic/room/\{roomId\}} for
\code{PARTICIPANT\_JOINED}/\code{LEFT}, \code{ROOM\_UPDATED}, and
\code{ROOM\_CLOSED}.

\paragraph{Agora token minting.} The service builds an Agora AccessToken~v2
by hand: it packs the app id, channel, uid, privileges (join, publish audio,
publish video, data) and expiry, signs the message with HMAC-SHA256 using the app
certificate, and zlib-compresses the result into the familiar ``\code{007}\ldots''
token. The media then flows client-to-Agora directly; the backend never sees it.

\paragraph{Proximity-voice engine.} The most novel piece of the service turns
avatar positions into voice groups deterministically. On each position snapshot
from the game server (pushed to the internal \route{/positions} endpoint), the
engine assigns any avatar standing inside a zone that has a \code{voiceRoomId} to
that zone's private channel, and clusters the remaining open-area avatars by a
union-find over mutual hearing radius; a cluster of two or more becomes an ad-hoc
channel named \code{prox-<workspaceId>-<smallestUserId>}, while a lone avatar
gets no channel. Only avatars whose group changed are notified, via a
\code{VOICE\_GROUP\_CHANGED} broadcast, and the endpoint returns the mapping of
each user to their Agora channel. This is the \emph{signalling} layer that decides
\emph{who should hear whom}; Agora carries the actual audio.

\paragraph{Events produced.} On room create/delete and membership change, a
\code{RoomChannelEvent} is published to the \code{room.exchange}; the chat service
consumes it to materialise or update the bound \code{ROOM} channel. Room creation
even \emph{rolls back} if the event cannot be published, so a room never exists
without its channel eventually appearing.

\section{Notification Service}
\label{sec:be-notif}

\paragraph{Responsibility.} The notification service (port 8082) has two jobs:
sending transactional \textbf{e-mails} (sign-up, login, OTP, password reset) from
versioned HTML templates, and delivering \textbf{in-app notifications} (currently
\code{TASK\_ASSIGNED}) by persisting them and pushing them over WebSocket. It is a
pure consumer of events and never publishes.

\paragraph{Stack.} Spring Web, Data Mongo, Data Redis, AMQP, Mail, and WebSocket.
An in-app \code{Notification} document records \code{userId}, \code{type},
\code{title}, \code{body}, an arbitrary \code{data} map, \code{read} state, and a
unique \code{eventId} for idempotency; e-mail events are not persisted.

\paragraph{Delivery.} A single \code{@RabbitListener} on
\code{notifications.queue} dispatches on the event type. E-mail events take an
idempotency claim in Redis (\code{SETNX notif:processed:\{eventId\}}, 24-hour
TTL), render the matching HTML template by \code{\{\{placeholder\}\}}
substitution, and send via SMTP --- releasing the claim on failure so a retry can
re-attempt. The \code{TASK\_ASSIGNED} event instead builds a notification, saves
it to MongoDB (the unique \code{eventId} making the insert idempotent), and pushes
a \code{\{action:"NEW\_NOTIFICATION", ...\}} envelope to the recipient's private
queue via Spring's user-destination routing. Consumption retries five times with
exponential backoff and then drops the message (there is no dead-letter queue in
this version).

\paragraph{Inbox API.} The web client reads history through
\route{GET /api/notifications} (paged), an \route{/unread-count}, and mark-read
and delete operations, plus the usual \route{/ws-ticket} for the socket. Because
offline users cannot receive a push, the persisted history backfills them on their
next inbox load --- push is best-effort, storage is the source of truth.

\section{The Unimplemented Modules}
\label{sec:be-stubs}

For an accurate account, three modules must be described as the placeholders they
are.

\begin{description}
  \item[Calendar service (port 8083).] A Maven scaffold containing only a POM,
  a minimal \code{application.properties}, and a README describing intended
  features (event scheduling, meeting invites, calendar sync). It has no Java
  source and is not wired into the gateway.
  \item[Tasks service (port 8085).] Likewise a scaffold. It is the intended
  future \emph{producer} of the \code{TASK\_ASSIGNED} events that the notification
  service already knows how to consume --- so the consumer side of task
  notifications exists and waits for this producer.
  \item[Shared library.] Declared as a module and intended to hold common DTOs,
  utilities, and exception handling, but currently containing no Java classes.
  Each service therefore defines its own exception handler, user-context helper,
  and DTOs independently.
\end{description}

These gaps are the clearest signposts for future work (Chapter~\ref{ch:conclusion}).

\section{Cross-service Consistency in Practice}

It is worth closing the chapter by watching the patterns of
Chapter~\ref{ch:architecture} act together in a single scenario:
\emph{a user accepts a workspace invitation}.

\begin{enumerate}
  \item The invitee opens the tokenised accept link; the request reaches the
  workspace service through the gateway with a validated \code{X-User-Id}.
  \item The workspace service marks the invitation \code{ACCEPTED} and creates
  (or reactivates) the invitee's \code{Desk} in one transaction.
  \item \emph{After} that transaction commits, it publishes a
  \code{WorkspaceChannelEvent} (member added) to RabbitMQ and returns success ---
  the user's request does not wait for anything downstream.
  \item The chat service, subscribed to that exchange, consumes the event
  (idempotently on \code{eventId}) and adds the user to the workspace's canonical
  channel, so their new channel appears a moment later.
\end{enumerate}

No service reached into another's database; identity was validated once; the
cross-service change was asynchronous, durable, and idempotent; and the user
experienced an immediate response with the rest catching up behind it. That
combination --- edge-validated identity, per-service data ownership, and
event-driven eventual consistency --- is the through-line of the entire backend.
